% Author: Simon-Pierre Boucher — contact@spboucher.ai
%
% ═══════════════════════════════════════════════════════════════════════
% 8. LIMITATIONS
% ═══════════════════════════════════════════════════════════════════════
\section{Limitations}
\label{sec:limitations}

We summarize the principal limitations in a structured format.

\begin{enumerate}[nosep]
    \item \textbf{Listing prices, not transaction prices.} The dependent variable is the asking price, which may differ from the realized sale price due to strategic pricing, negotiation, and market conditions. Listing-price gradients are not necessarily equivalent to transaction-price implicit prices.

    \item \textbf{No causal identification.} All estimates are conditional associations. Without exogenous variation, we cannot distinguish the causal effects of attributes from sorting, supply constraints, and omitted variables \citep{ekeland2004identification, roberts2013endogeneity}.

    \item \textbf{Spatial dependence not modeled.} Moran's $I = 0.27$ confirms substantial spatial autocorrelation in OLS residuals. The paper does not estimate spatial lag, spatial error, or geographically weighted regression models, which may affect both efficiency and consistency of estimates.

    \item \textbf{Random validation overstates ML performance.} The random train-test split allows spatial leakage. Under geographic holdout, XGBoost $R^2$ drops from 0.833 to 0.425--0.547 depending on specification.

    \item \textbf{Standardization complicates interpretation.} The standardized OLS coefficients are not directly interpretable as elasticities or semi-elasticities. We provide an unstandardized specification for economic interpretation but note that some readers may conflate the two.

    \item \textbf{Missing data and imputation.} Year built (19.4\%), lot size (16.3\%), and transit score (70.1\%) have substantial missingness. State-median imputation preserves geographic variation but attenuates the lot-size coefficient by a factor of three.

    \item \textbf{Climate risk variables entirely missing.} Despite their growing importance, flood, fire, heat, wind, and air risk factors could not be incorporated \citep{baldauf2020does}.

    \item \textbf{SHAP values are not structural implicit prices.} SHAP provides prediction decompositions, not marginal willingness-to-pay estimates. Cross-model stability supports the ranking but not the economic interpretation.

    \item \textbf{Quantile regression subsample.} Quantile regression is estimated on 150,000 observations (19\% of the full sample) for computational reasons. Most coefficients are stable (CV $< 8\%$) but age is not (CV = 78.6\%).

    \item \textbf{Non-representative sample.} The Zillow listing dataset overrepresents Southern states, active listings, and possibly certain property types. Results may not generalize to the full U.S.\ housing stock.

    \item \textbf{Data quality concerns.} Bike Score exceeds 100 for 38 observations; some lot sizes are implausibly large; parking-space data are noisy. These affect a small fraction of observations but may influence tail estimates.

    \item \textbf{No hyperparameter tuning via cross-validation.} XGBoost and LightGBM hyperparameters were set based on common defaults rather than optimized via nested cross-validation. Performance could potentially improve with systematic tuning.

    \item \textbf{Geographic holdout design is conservative.} Holding out 10 entire states is a stringent test; finer geographic blocking (county-level or MSA-level) might yield intermediate performance estimates that are more relevant for some AVM applications.

    \item \textbf{Single cross-section.} The data represent a snapshot of active listings at one point in time. Temporal variation in hedonic gradients---due to market cycles, interest rate changes, or policy shifts---cannot be examined.
\end{enumerate}
